\documentclass[UTF8,a4paper,zihao=-4,fontset=fandol]{ctexart}
% Shared formatting for the course report and research manuscript.
\usepackage[margin=2.2cm,headheight=14pt]{geometry}
\usepackage{amsmath,amssymb,bm,booktabs,tabularx,longtable,array}
\usepackage{graphicx,xcolor,tikz}
\usetikzlibrary{arrows.meta,positioning,calc}
\usepackage{caption,float,placeins}
\usepackage[hidelinks,unicode]{hyperref}
\usepackage{url}
\definecolor{signalblue}{HTML}{205A87}
\definecolor{referenceorange}{HTML}{A95822}
\definecolor{softgray}{HTML}{F0F3F5}
\captionsetup{font=small,labelfont=bf,labelsep=quad}
\renewcommand{\arraystretch}{1.2}
\setlength{\tabcolsep}{4.5pt}
\setlength{\parskip}{0.14em}
\setlength{\emergencystretch}{3em}
\linespread{1.22}
\pagestyle{plain}
\setcounter{tocdepth}{2}
\newcommand{\um}{\,\mu\mathrm{m}}
\newcommand{\ns}{\,\mathrm{ns}}
\newcommand{\mV}{\,\mathrm{mV}}
\newcommand{\fC}{\,\mathrm{fC}}
\newcommand{\V}{\,\mathrm{V}}
\newcommand{\sourcepath}[1]{{\footnotesize\nolinkurl{#1}}}
\graphicspath{{figures/}}
\hypersetup{pdfauthor={},pdfsubject={可追溯的阶段研究成果；独立测试与工艺扩展尚待最终验收}}

% Fixed local-ledger snapshot; do not interpret RUNNING as a live process check.
\newcommand{\EvidenceDate}{2026-10-11 05:02\,UTC$-08{:}00$}
\newcommand{\TrainVerified}{64}
\newcommand{\HoldoutVerified}{19}
\newcommand{\HoldoutRunning}{1}
\newcommand{\HoldoutMissing}{12}
\newcommand{\ExtrapVerified}{0}
\newcommand{\ExtrapRunning}{0}
\newcommand{\ExtrapMissing}{32}

\hypersetup{pdftitle={AI辅助下的多特征尺寸DRAM单粒子效应评估：阶段成果整理稿}}
\begin{document}
\begin{titlepage}
\centering
\vspace*{1.3cm}
{\zihao{2}\bfseries 项目设计课程报告\par}
\vspace{1.8cm}
{\zihao{1}\bfseries AI辅助下的多特征尺寸DRAM\\[0.3em]单粒子效应评估\par}
\vspace{0.8cm}
{\zihao{3}\color{signalblue}阶段成果整理稿\par}
\vspace{2cm}
\begin{tabular}{rl}
指导教师：&李洋\\[1em]
项目成员：&\rule{8cm}{0.4pt}\\[1em]
院系及班级：&\rule{8cm}{0.4pt}\\[1em]
证据快照：&\EvidenceDate
\end{tabular}
\vfill
\begin{minipage}{0.85\textwidth}
\small
本稿整理已留存的真实仿真、模型输入、数据处理与验证证据。当前版本为
R50-main阶段快照，训练已完成、测试预测已冻结，尚未完成全部测试评分和扩展验证。
成员身份、实际分工、教师审核及答辩记录由参与者按真实情况补充。
\end{minipage}
\vspace{1cm}
\end{titlepage}

\pagenumbering{Roman}
\begin{abstract}
围绕课程提出的AI辅助多特征尺寸存储器辐射效应评估任务，本项目选用简化
1T1C DRAM单元，建立Windows本机与Linux虚拟机之间的远程访问式仿真流程。
限定中文模板用于表达并核对参数矩阵，Geant4计算硅内粒子的电离沉积，
经能量、电荷与径迹长度换算形成TCAD等效源，Sentaurus完成器件和电路瞬态响应求解。
研究依次经历单事件链路贯通、受控初态和读出判据建立、参考偏移诊断及统一
50~$\Omega$预充版本重建。当前R50-main已有32个独立训练事件及64个配对位响应，
并在测试响应可见前冻结64条逐测试位预测。训练工况分组交叉验证显示，四个岭回归
模型的电荷和裕量误差均低于训练均值基准；这属于训练阶段评价，域内留出和尺寸
外推的最终精度仍待完整仿真评分。报告重点说明模型假设、参数传递、匹配参考、
独立事件分母及数值验证，并保留预充失败与物理修订的证据。结果支持一条可追溯的
教学研究流程，尚不构成真实商业DRAM位翻转率、辐射截面或工艺校准结论。

\noindent\textbf{关键词：}DRAM；单粒子效应；Geant4；TCAD；AI辅助；岭回归
\end{abstract}
\clearpage
{\small\linespread{1.0}\selectfont\tableofcontents}
\clearpage
\pagenumbering{arabic}
\ctexset{section={name={第,章},number=\arabic{section}}}
\section{任务目标、基本原理与可行性}
\label{cr:objectives}
\subsection{课程任务与研究对象}
课程要求理解DRAM、单粒子效应和AI基本方法，在个人计算机与远程计算资源之间
搭建仿真环境，完成基础模型、自然语言改参、Geant4到TCAD传递、多参数扫描以及
拟合预测验证，并由2至3名学生合作完成设计、演示和总结\cite{course}。
据此，项目将目标分解为三个可检查的问题：能否把一个明确的粒子条件转换为
可审查的器件源项；能否在相同电学前提下比较不同尺寸和粒子条件的响应；能否
用部分已求解工况预测尚未参与拟合的工况，再由真实TCAD结果检验预测。

本报告选择1T1C DRAM作为对象。一个访问晶体管控制存储节点与位线的连接，
存储电容上的电荷表征信息。保持时访问管关断，读出时字线打开访问管，
存储电容与预充位线发生电荷共享，感测电路据微小差分判别位值\cite{keeth2001}。
因此，“粒子造成多少电荷变化”和“读出是否仍有足够裕量”是相关而不同的问题。
项目既提取保持段电荷扰动，也建立读出指标，避免仅凭存储节点电压下降量判断
完整存储单元的逻辑状态。

\subsection{单粒子效应与数值描述}
载能粒子穿过半导体时产生电子空穴对，载流子在内部电场中漂移，并发生扩散、
复合及端口收集。当扰动到达敏感节点，可能形成瞬态电流或改变储存电荷，
进一步影响电路判读\cite{label1996}。单粒子瞬态、判据上的错误读出与锁存后
持续存在的位翻转并不等价。本项目尚无真实感放锁存和恢复电路，故主要研究
电荷扰动及明确采样时刻的判据裕量。

在漂移扩散框架中，粒子作用以产生率$G$进入电子、空穴连续性方程：
\begin{equation}
\frac{\partial n}{\partial t}=\frac{1}{q}\nabla\!\cdot\mathbf{J}_n+G-R,\qquad
\frac{\partial p}{\partial t}=-\frac{1}{q}\nabla\!\cdot\mathbf{J}_p+G-R.
\end{equation}
这些方程与泊松方程及外部电路方程耦合，决定存储节点和位线的瞬态变化。
Geant4负责粒子输运及能量沉积，TCAD负责沉积之后的半导体电学响应。
二者之间的接口必须明确材料、坐标、单位、源项形状及有效厚度；工具同时运行
并不足以证明它们描述了相同的物理对象。

\subsection{可行性与实际设计过程}
实际环境由Windows本机承担请求解析、批次管理、归档和后处理，由已有Linux
虚拟机运行Geant4与Sentaurus。本机通过严格SSH入口访问虚拟机，保留命令、
返回状态和输入输出摘要。这实现了远程访问式工作流，但现有材料没有提供
实验室独立服务器或正式资源交接记录，不能把虚拟机使用记录替代课程要求中的
资源交接。现有软件、许可证和脚本已能支持单核串行求解，主要成本来自TCAD
瞬态和必要复验，而非表格或绘图。

设计采用逐级建立证据的方式。v1先用小矩阵贯通真实输运、器件响应与外推；
v2加入统一初态、0/1配对和读出判据；D1进一步追查理想参考留下的读出偏移；
随后R50-main以匹配中性参考及统一预充版本建立更完整的矩阵。
这种过程让前一阶段暴露的问题成为后一阶段的设计依据。
当前报告以共享证据快照为截止点，完成状态见表\ref{cr:progress}；
历史阶段保留其原条件，不用新模型重写旧结论\cite{project}。

\section{基础模型与参数来源}
\label{cr:model}
\subsection{几何、电路与物理假设}
基础模型为自编二维平面NMOS加理想存储电容的1T1C单元。设栅长为$L$，
硅体宽度为$L+0.6\um$、深度为$1.2\um$，二维有效第三维厚度为$1\um$。
主模型氧化层厚度为$5\,\mathrm{nm}$，体区为P型，体掺杂浓度
$10^{17}\,\mathrm{cm}^{-3}$，源漏高斯峰值为$10^{20}\,\mathrm{cm}^{-3}$，
结深参数为$0.12\um$。栅功函数设为$4.6\,\mathrm{eV}$，
温度为$300\,\mathrm{K}$，外接存储电容为$10\,\mathrm{fF}$。
比较的四个栅长是$0.18/0.25/0.35/0.45\um$，表示访问晶体管的几何长度，
并不表示商业DRAM工艺节点。

模型使用Fermi统计、OldSlotboom带隙修正、掺杂和电场相关迁移率以及
SRH、Auger复合。器件采用漂移扩散求解，电路节点同时参与瞬态计算。
选用这种结构的依据是它保留了访问管、电荷存储及电荷共享三项基本机制，
输入规模适合当前计算环境，也便于检查粒子源与电路响应之间的关系。
真实DRAM可能采用复杂沟道、埋栅和立体电容，简化结构不能替代这些工艺的
电学标定，更不能由栅长数值推断产品代际。

\begin{table}[htbp]
\centering\small
\caption{主模型参数与证据性质。数值为本项目冻结的探索参数，不能统一称为文献实测值。}
\label{cr:parameters}
\begin{tabularx}{\textwidth}{lllX}
\toprule
参数&主模型取值&用途&来源或适用范围\\
\midrule
栅长$L$&0.18至0.45 $\mu$m&几何扫描&自编平面访问管，四点离散取值\\
氧化层&5 nm&主矩阵&4/6 nm为后续受控变体，尚无最终结果\\
存储电容&10 fF&储存电荷&1T1C原理与已有混合模式示例支持形式，未校准器件\\
位线与参考电容&100 fF&读出电荷共享&统一电路假设，便于匹配比较\\
初态SN&0.1/0.9 V&存储0/1&理想钳位初始化，非真实写入校准\\
中性参考SN&0.5 V&消除共模偏移&同器件、同网格与时序的参考支路\\
预充导通电阻&50 $\Omega$&R50-main读出&由独立probe支持的物理修订\\
\bottomrule
\end{tabularx}
\end{table}

\subsection{从参数来源到可检验假设}
参数来源分为原理、工具约定、已有模型形式和本项目假设四类。
DRAM教材支持访问晶体管与电容的基本工作机制；Geant4资料支持输运过程和
电磁模型的含义\cite{keeth2001,g4em}。已有混合模式DRAM示例说明外接理想
电容的建模方式可行，但本项目没有复制并校准其完整工艺工程。
几何尺寸、掺杂、位线电容及裕量门槛仍按探索假设处理。用明确的来源分层，
可以把“软件支持这个模型”与“模型代表某种真实器件”区分开来。

R50-main的0/1状态通过理想SN钳位获得，钳位在$0.6\ns$释放，
并在$0.95\ns$检查初态误差不超过$1\mV$。这一设置减少不同栅长写入程度
不同所带来的混杂，使比较更接近“相同初态下的访问管几何变化”。
然而它只控制初始电压，不证明实际写入时间、能耗或访问管写入能力。
v1曾采用实际访问管充电，不同栅长的粒子前电压不同；
因此两版数据适合解释设计改进，不能拼接成一条统一的尺寸响应曲线。

\subsection{保持、预充与匹配参考}
粒子在$1\ns$注入，保持指标在$2.8\ns$提取。位线随后经预充支路升至
$0.5\V$，在$3\ns$释放；字线在$3.05$至$3.10\ns$升至$1.8\V$，
读出在$4\ns$采样。R50-main使用$50\,\Omega$导通电阻，
关断电阻为$10^{20}\,\Omega$，不把理想开关的电压跳变当作预充完成。

匹配中性参考在相同尺寸、网格、偏压及时序下取SN=$0.5\V$。
批量计算使用信号位线与中性参考位线之差，并先通过真实双支路仿真核对。
这种参考随访问管打开而产生相应的耦合和电荷交换，比统一减去一个常数更适合
处理状态相关偏移。参考中不注入粒子；物理双支路中采用独立器件定义，
以防同一HeavyIon设置同时作用于信号与参考。

\section{AI交互平台与Geant4到TCAD的传递}
\label{cr:platform}
\subsection{自然语言、结构化设计与执行边界}
AI在项目中辅助阅读文献、组织假设、生成和检查脚本、定位失败、提取结果及
撰写解释。可以被复验的交互接口则是限定中文模板：请求明确列出栅长、
粒子、能量、方向、重复数和范围，解析后得到结构化工况清单。
每个工况包含几何条件、粒子条件、种子、配对位值及训练或测试身份。
接口拒绝模糊尺寸、未支持角度和超出冻结范围的要求，避免把自然语言中的
省略或歧义直接传给求解器。

v1已有中文请求生成完整矩阵并进入真实执行链路的记录。
当前R50演示入口提供\texttt{preview}、\texttt{audit}和\texttt{locate}：
分别预览映射、检查接受与拒绝行为、定位指定事件的证据路径。
它只说明请求与已冻结R50设计如何对应，不自动提交新求解。
$50\,\Omega$来自版本化物理修订，而非解析器从句子中猜测出的参数。
自然语言成功解析是软件层的检查，实际模型正确性仍由后续物理与数值证据判定。

\begin{figure}[htbp]
\centering
% Conceptual workflow, not a measured result or screenshot.
\resizebox{\linewidth}{!}{%
\begin{tikzpicture}[
 block/.style={draw=signalblue,rounded corners=2pt,fill=softgray,
 text width=4.2cm,minimum height=1.5cm,align=center,font=\small},
 arrow/.style={-{Latex[length=2.4mm]},thick,draw=signalblue},
 node distance=0.7cm and 0.6cm]
\node[block] (request) {限定中文请求\\尺寸、粒子、角度\\独立事件数};
\node[block,right=of request] (matrix) {受约束的工况清单\\单位、范围、版本与种子};
\node[block,right=of matrix] (transport) {Geant4 硅内电离沉积\\$E_{\rm ion}\ \rightarrow\ Q_{\rm dep}$\\有限径迹及单位检查};
\node[block,below=of transport] (tcad) {TCAD 器件与电路求解\\等效径迹与受控初态\\读出瞬态与端口检查};
\node[block,left=of tcad] (metric) {匹配中性参考与同位对照\\$D,\ M,\ \Delta Q,\ \Delta M$};
\node[block,left=of metric] (model) {按工况分组的响应拟合\\冻结预测后检验未见响应};
\draw[arrow] (request)--(matrix);
\draw[arrow] (matrix)--(transport);
\draw[arrow] (transport)--(tcad);
\draw[arrow] (tcad)--(metric);
\draw[arrow] (metric)--(model);
\node[below=0.40cm of metric,align=center,font=\small,text width=14.5cm]
 {每一层绑定输入版本、运行记录和来源哈希；科学质量门槛与执行状态分别记录。};
\end{tikzpicture}}

\caption{AI辅助联合评估的工作流程示意。自然语言请求经结构化审查与冻结后，
连接粒子输运、等效源传递、器件求解和独立评价；此图为流程示意，不代表未完成阶段已有结果。}
\label{cr:workflow}
\end{figure}

\subsection{输运记录与等效源转换}
Geant4在硅与真空组成的简化几何中运行电磁option4及步长限制，
记录每一步的位置、时间、总沉积能、非电离沉积能和轨迹身份。
主矩阵采用$5\,\mathrm{MeV}$的$\alpha$粒子与质子，方向为0°和30°，
表面入射位置与各栅长模型保持约定关系。30°径迹可能从侧面退出，
因此不能机械地把硅深度除以角度余弦视为所有工况的有效径迹长度。
实际几何交点及源文件共同决定传递参数。

只把硅内电离能量计入电子空穴对产生。设逐步电离沉积总量为
$E_{\mathrm{ion}}$，使用硅中平均$3.6\,\mathrm{eV}$产生一对载流子的近似，
得到
\begin{equation}
N_{\mathrm{eh}}=\frac{10^6E_{\mathrm{ion}}[\mathrm{MeV}]}{3.6},\qquad
Q_{\mathrm{dep}}=qN_{\mathrm{eh}},\qquad
\mathrm{LET}_{\mathrm{eq}}=\frac{Q_{\mathrm{dep}}}{\ell}.
\label{cr:transfer}
\end{equation}
式中$\ell$是约定的硅内等效直径迹长度。TCAD以半径$0.03\um$的
Gaussian HeavyIon源注入等效电荷，并使用一致的二维有效厚度和电荷单位。
这一方法保留事件的电离总量，便于小规模参数扫描；逐步输运CSV仍单独保留，
以便识别总量近似忽略的纵向、横向和时间分布信息。

已有TCAD与Geant4的SRAM评估研究提供了联合方法的先例\cite{chen2011}；
硅二极管中重离子、聚焦X射线和脉冲激光的瞬态比较，则提示不同电荷产生方式
需要分别处理\cite{ryder2021}。这些文献不能直接校准本项目的DRAM模型。
本项目采用总量等效直径迹而非完整三维产生率网格，
故需要分别检查转换算术与TCAD实际源积分。前者核对单位和总电荷，
后者检查离散网格、径迹截断和时间积分之后是否仍近似注入目标电荷。
二者通过并不意味着空间分布误差消失。

\subsection{可追溯执行与数据流}
执行流程将设计、输入deck、Geant4事件三文件、求解日志及输出归档绑定到唯一
run ID，并用SHA-256检查内容一致性。续算只允许使用同几何、同电路、同网格和
已核对合同的保存点；改变网格或参考结构后重新建立相应初态。
主运行器负责归档，分析器从固定原始输出提取曲线，图表再引用这些曲线和摘要，
形成“请求—输入—源—求解—指标—图”的连续来源关系。

这种设计尤其适合AI辅助操作。AI能够快速生成大量命令，也可能误读旧摘要、
把同名文件当作同一内容，或把后台退出当成远端求解停止。
因此流程把可机器检查的版本与完成条件写入入口，而不是依赖对话记忆。
报告中的数据图属于原始数值后处理；真实Sentaurus窗口截图和原生场图另行标注，
防止用美观的曲线图替代软件求解证据。

\section{实验矩阵、预测设计与验证方法}
\label{cr:design}
\subsection{固定矩阵与响应定义}
R50-main主设计包含四种栅长、两种粒子和两个方向，每个组合预定四个独立种子，
共64个独立粒子事件；每个事件分别施加于存储0和存储1，形成128个响应槽位。
配对设计有助于识别位值不对称，但同一沉积事件的两种初态不构成两个独立入射。
数值复跑、网格细化、参考诊断和氧化层配对同样不增加独立事件数。
所有种子与工况在读取正式响应前登记，失败不通过更换种子补成有利样本。

R50-main使用如下连续指标：
\begin{align}
\Delta Q&=C_{\mathrm{store}}\,[V_{\mathrm{SN,event}}(2.8\ns)
 -V_{\mathrm{SN,control}}(2.8\ns)],\\
D(t)&=V_{\mathrm{BL,signal}}(t)-V_{\mathrm{BL,neutral}}(t),\\
M&=(2b-1)D(4\ns),\qquad
\Delta M=M_{\mathrm{control}}-M_{\mathrm{event}}.
\label{cr:metrics}
\end{align}
其中control是同位值无粒子对照，neutral是中性参考，两者功能不同。
$\Delta Q$保留符号，负值表示事件SN低于对照；$\Delta M>0$表示裕量退化。
对存储0而言，节点进一步下降可能使判据裕量增加，此时$\Delta M$可以为负，
不应把所有变化都强制写成非负的“损伤”。

主判据取$\varepsilon=5\mV$：$M>\varepsilon$为正确，
$M<-\varepsilon$为判据误读，其余为临界。
同时保留$2.5/10\mV$阈值结果。接近边界而未完成规定数值复验、
或粗细离散分类不一致时，另列数值不确定；失败和缺失保持各自状态。
统计按冻结分母报告，不用可判子集替代预定样本。

\subsection{分组拟合与测试前冻结}
训练集取$L=0.18/0.35\um$，共32个事件、64个响应；
$0.25\um$为域内留出，$0.45\um$为尺寸外推，各16个事件、32个响应。
同一物理事件及其配对结果保持同组，避免相同源信息同时出现在训练和测试中。
预测使用真实Geant4沉积量，目标是在给定沉积与无粒子条件下近似TCAD响应。
它有望减少后续器件求解成本，但仍须获得粒子输运输入。

分别针对两种位值的$\Delta Q$与$\Delta M$建立四个岭回归模型，
特征为$Q_{\mathrm{dep}}$、$Q_{\mathrm{dep}}/L$和
$Q_{\mathrm{dep}}\sin\theta$。这些项以电荷注入规模、尺寸及方向为物理启发，
并不是已推导出的通用器件定律。与训练均值基准比较，可检查回归是否比
忽略输入变量更有信息；数据驱动单粒子响应建模的相关研究也为这种代理思路
提供背景\cite{xu2021}。

训练采用留一物理工况交叉验证（LOCO）。训练区有八个尺寸、粒子、角度组合，
每折将一个完整组合的四个重复留出；两位保持相同的分组边界。
标准化均值和方差只由本折训练数据计算。候选正则化参数为
$0,0.01,0.1,1,10$，以八折RMSE的算术均值选择，不由测试集选模。
岭目标为
\begin{equation}
\min_{\beta_0,\boldsymbol\beta}
\sum_i(y_i-\beta_0-\boldsymbol z_i^{\mathsf T}\boldsymbol\beta)^2
+\lambda\|\boldsymbol\beta\|_2^2,
\end{equation}
截距不惩罚。最终参数用完整训练集拟合，在首次测试TCAD提交前冻结模型、
标准化参数、Geant4测试输入、64条逐位预测及其哈希。
留出和外推完成后仅计算误差，不再按测试误差修改模型。

\subsection{分层验证与工程验收}
每个有效响应需满足完整求解到$4\ns$、有限数据、初态误差不超过$1\mV$、
预充误差不超过$10\,\mu\mathrm{V}$，并检查逐相KCL、四端口闭合和
后向欧拉右端点电荷积分。粒子源实际积分与目标Q的偏差门槛为5\%，
参考支路源项应为零。对物理双支路和离线参考同时比较两侧BL及差分，
防止两侧误差偶然抵消。

预登记代表点分别进行半时间步和0.7网格尺度检查，并为每个变体匹配自己的
对照和中性参考。$\Delta M$变化门槛为$0.5\mV$、$\Delta Q$为$0.01\fC$，
同时核对分类稳定性。训练交叉验证、数值细化和独立测试解决不同问题：
前者检验训练工况内的拟合选择，中者检查离散敏感性，后者检验未见尺寸条件。
任何一项通过都不能自动替代另外两项。

\section{已完成结果与验证证据}
\label{cr:results}
\subsection{证据快照与阶段进度}
表\ref{cr:progress}取自本报告共享快照，截止时间为\EvidenceDate。
R50-main的训练响应已完成\TrainVerified{}项；留出已验证\HoldoutVerified{}项，
账本在途\HoldoutRunning{}项，尚缺\HoldoutMissing{}项；外推已验证
\ExtrapVerified{}项，账本在途\ExtrapRunning{}项，尚缺\ExtrapMissing{}项。
这些计数描述逻辑响应槽位；在途是快照记录，不表示报告生成时重新核查过实时进程。
工程重试保留原记录但不增加独立样本分母。
本稿不把正在运行的输出纳入评分，也不因后续文件变化自动改写本次快照。

\begin{table}[htbp]
\centering\small
\caption{R50-main在共享证据快照时的完成状态。训练与测试预测已冻结，
未完成的求解及后续验证均按真实状态列出。}
\label{cr:progress}
\begin{tabular}{llrrrr}
\toprule
阶段 & 栅长 / $\mu$m & 计划响应 & 已验证 & 在途 & 待提交 \\
\midrule
训练 & 0.18 / 0.35 & 64 & 64 & 0 & 0 \\
域内留出 & 0.25 & 32 & 19 & 1 & 12 \\
域外外推 & 0.45 & 32 & 0 & 0 & 32 \\
\bottomrule
\end{tabular}

\end{table}

\subsection{历史链路验证与读出偏移诊断}
v1完成16个固定事件工况，覆盖四栅长、两粒子和两方向。
匹配对照修正的SN电压下降范围为$0.018553$至$0.490764\V$；
源项积分误差为0.2036\%至1.2422\%。
其8点训练、4点域内留出、4点域外评价中，域内MAE为$4.392\mV$，
域外MAE为$3.414\mV$。域外预测有先于参考TCAD完成的登记记录，
而域内响应在系数文件冻结前可见，所以只称留出评价。
这批结果证明了小矩阵的真实贯通和两阶段代理能力，尚无独立事件统计。

v2用受控初态建立读出流程，在$L=0.25\um$、$5\,\mathrm{MeV}$、
垂直入射下完成每种粒子四个独立事件。原理想参考和$5\mV$门槛下，
$\alpha$存储1为1项正确、3项临界，其他位值组均正确。
明确误读为每种粒子0/4，并不能证明误读概率为零；完全可判且零事件的
单侧95\%上界在$N=4$时仍达52.7\%。这些历史分类保持原判据，不按后来的
参考方案追溯替换。

D1通过九个无粒子工况说明参考选择为何重要。中性SN=$0.5\V$、标准字线下，
理想参考保留了约$+12.370\mV$的读出差分，关闭字线后仅为
$-0.003816\mV$。当无粒子SN与原$\alpha$事件读前状态相匹配时，
理想参考得到$+2.454961\mV$；离线中性参考和实际匹配支路均得到
$-9.915036\mV$。端口积分及电容净电荷闭合支持字线耦合与访问管交换的解释。
这推动了R50-main的匹配参考设计，但D1不是新的粒子样本，也不是R50响应。

\begin{figure}[htbp]
\centering
\includegraphics[width=0.88\textwidth]{d1-reference-comparison.pdf}
\caption{历史D1无粒子诊断的参考方式比较。$L=0.25\um$，存储电容10 fF、
位线100 fF，信号SN匹配原$\alpha$事件的读前电压；读出在4 ns比较，差分单位mV。
理想、离线中性及物理匹配参考属于同一诊断比较，不对旧v2事件重新分类。}
\label{cr:d1}
\end{figure}

\subsection{预充修订与R50参考验证}
旧$100\,\Omega$版本的一个固定$\alpha$事件在$3\ns$的BL预充误差为
$-13.297362\,\mu\mathrm{V}$，超过原$10\,\mu\mathrm{V}$门槛。
随后独立$50\,\Omega$ probe的匹配事件误差为
$-4.731204\,\mu\mathrm{V}$，control与neutral亦通过；
源、初态和其余时序保持对应关系。
因此项目建立独立R50-main物理版本，从头生成所需基准与参考，
而未把原失败改判或把probe填入正式矩阵。

\begin{figure}[htbp]
\centering
\includegraphics[width=0.88\textwidth]{r50-precharge.pdf}
\caption{旧100 $\Omega$与独立50 $\Omega$ probe中同一固定事件的预充误差绝对值。
精确采样时刻3 ns，量为$|V_{\mathrm{BL}}-0.5\,\mathrm{V}|$，单位$\mu$V；
固定上限为10 $\mu$V。旧失败保留，probe不计入R50-main训练或测试样本。}
\label{cr:precharge}
\end{figure}

R50-main现有20项参考建立与验证证据：四尺寸的0/1无粒子对照及中性参考共12项，
以及四尺寸乘两位的物理双支路比较8项。
小批量粒子响应及两项源隔离检查也已完成；预登记代表点的匹配数值检查通过。
图\ref{cr:reference}展示这些基准的受控初态和代表尺寸的无粒子读出波形。
前置结果支撑当前训练数据的参考定义和接口使用，但不替代主矩阵完成后的
边界细化，也不证明所有尺寸、所有事件均达到空间网格独立。

\begin{figure}[htbp]
\centering
\includegraphics[width=0.88\textwidth]{r50-reference.pdf}
\caption{R50-main无粒子基准与中性参考。左图为四栅长、存储0/中性/存储1
共12项的粒子前SN电压，单位V；右图为$L=0.18\,\mu$m三种初态在3至4 ns
读出阶段的BL电压，单位V。全部来自50 $\Omega$参考快照，未混入粒子主响应或旧D1数据。}
\label{cr:reference}
\end{figure}

\subsection{训练响应与冻结预测}
R50-main已有32个独立训练事件、64个配对响应，见图\ref{cr:training}。
在$5\mV$主阈值下，64行训练记录中48行为正确、10行为判据误读、
4行为临界、2行为数值不确定。这个汇总用于描述训练数据构成，
不能把64行当成64次独立粒子入射，也不是持久SEU率。
连续目标与离散分类应同时保留：临界或近边界记录仍可能提供有用的连续响应信息，
但其类别判断需要相应数值证据。

\begin{figure}[htbp]
\centering
\includegraphics[width=0.94\textwidth]{r50-training-response.pdf}
\caption{R50-main已完成训练集的连续响应。栅长0.18/0.35 $\mu$m，
5 MeV $\alpha$/质子、0°/30°，共32个独立事件、64个配对位响应。
$\Delta Q$为2.8 ns同位无粒子对照修正电荷，单位fC；
$\Delta M$为4 ns匹配中性参考下的裕量损失，单位mV。}
\label{cr:training}
\end{figure}

四个位值与目标组合均选取$\lambda=0.1$。表\ref{cr:cv}列出训练LOCO结果：
两个电荷目标的RMSE为0.102888和0.112903 fC，
两个裕量目标为0.955160和1.012491 mV；训练均值基准分别为
2.279967、2.515710 fC和21.458870、23.383861 mV。
这说明选定特征在已采样训练工况中包含明显的响应信息，而不只是拟合一个常数。
模型与64条测试位预测已经冻结，冻结记录中测试响应可见数为零。
尚未完成的测试评分不能由较小的交叉验证误差代替。

\begin{table}[htbp]
\centering\small
\caption{R50-main训练工况LOCO评价。数值是八个留出工况折RMSE的算术均值，
按存储位和目标分别评价；不是独立测试集的误差。}
\label{cr:cv}
\begin{tabular}{clrr}
\toprule
存储位 & 目标（单位） & 岭回归 RMSE & 均值基准 RMSE \\
\midrule
0 & $\Delta Q$ (fC) & 0.102888 & 2.279967 \\
0 & $\Delta M$ (mV) & 0.955160 & 21.458870 \\
1 & $\Delta Q$ (fC) & 0.112903 & 2.515710 \\
1 & $\Delta M$ (mV) & 1.012491 & 23.383861 \\
\bottomrule
\end{tabular}

\end{table}
\begin{figure}[htbp]
\centering
\includegraphics[width=0.9\textwidth]{r50-training-cv.pdf}
\caption{R50-main训练LOCO中岭回归与训练均值基准的比较。
电荷目标以fC、裕量目标以mV表示；四个模型均仅使用训练数据选取$\lambda=0.1$。
该图不表示域内留出或尺寸外推已经验证。}
\label{cr:cvplot}
\end{figure}
\FloatBarrier

\section{失败分析、结果解释与适用局限}
\label{cr:discussion}
\subsection{工程完成与科学有效性的区别}
预充失败表明，完成归档、求解器正常结束和读出分类为正确是三个不同判断。
一个运行即使到达$4\ns$，仍可能未满足预先规定的电学条件。
旧$100\,\Omega$事件的读出裕量远离错误边界，但其预充状态超差，
因此仍不能作为有效主响应。R50修订改变了外部电路，而不是更改误差容限；
它需要独立版本、重新初始化和新的匹配参考，不能称为单纯数值收敛改进。

历史过程还出现输入名称导致的解析错误、数据收集中断和工程重试。
处理原则是保留原输入、日志和失败身份，先确认远端实际状态，再恢复收集或
进行获准的唯一重试。成功重试只补足原逻辑槽位，不增加粒子事件数。
这类证据适合写入附录以说明流程可靠性；正文应着重解释它对数据纳入和结论的影响，
避免用大量命令输出取代研究论证。

\subsection{电荷变化、参考响应与分类的关系}
存储节点在粒子作用后降低，不意味着存储0和存储1的读出裕量都会变差。
式\eqref{cr:metrics}对位值引入符号后，同方向的电压扰动可能削弱存储1，
却增大存储0的裕量。源注入总电荷也不等于存储电容的净变化：
载流子可被不同端口收集、复合，并通过寄生电容与外电路交换电荷。
因此，不能以$Q_{\mathrm{dep}}/C$直接替代真实瞬态求解，也不应把固定电容下
电压和电荷的等价换算当作两项独立验证。

D1进一步说明，参考支路决定了哪些共模响应会留在差分中。
理想固定参考不会跟随访问管的读出耦合，而匹配中性支路会响应同样的字线操作。
但匹配并非完美消除所有状态依赖：信号与参考的初态不同，其沟道和结区演化
也可能不同。R50采用物理双支路一致性检查，是为了确认离线差分的实现，
并未因此完成真实感放失调、噪声和锁存恢复的校准。

\subsection{预测效果的含义与潜在偏差}
训练交叉验证明显优于均值基准，支持继续检验冻结代理模型；
它尚不能证明模型在$0.45\um$或新能量条件下仍有相同精度。
特征包含已知Geant4沉积量，使模型获得粒子与路径差异的重要信息，
也限定了预测的用途：它替代的是后续TCAD响应近似，不是完整输运过程。
当前仅有两个粒子类型、一个能量和两个方向，不能据此讨论连续空间辐射谱。

分组交叉验证减少了同工况重复跨折导致的乐观评价，但样本仍来自共同的简化模型。
未校准几何、理想电容、受控初态和统一材料假设可能共同形成稳定的系统偏差，
这种偏差不会因训练误差变小而被发现。
后续测试应同时报告逐点残差、分层误差及基准比较；若外推误差明显增大，
应分析特征与物理状态的适用边界，而非通过回看测试集调整模型后仍称原测试有效。

\subsection{仍需完成的技术与课程工作}
当前优先完成已冻结测试响应的真实求解和评分，再执行既定4/6 nm氧化层变体、
主矩阵最近边界点的匹配细化，以及R50原生场图的实际视觉验收。
氧化层扩展只讨论固定硅入口事件下的TCAD厚度敏感性；
Geant4仍为硅/真空输运，不能把它写成完整材料栈工艺变化。
电磁option4未包含核反应物理，质子结果不能解释为已验证的核反冲SEE。

若进一步面向器件研究，需要获得具有出处的DRAM结构和电学参考，
校准访问管、写入保持、感放与恢复，并按应用目标扩展粒子能量、位置与统计规模。
若面向课程验收，还需学生亲自解释模型、运行限定演示、完成真实分工与教师审核。
AI生成脚本和报告可以支撑这些活动，但不能替代参与人数、学生理解或考核事实。
现阶段可交付的是已验证技术产物及其明确边界，最终课程结论应由完整技术与教学
证据共同支持。

\section{阶段结论}
\label{cr:conclusion}
本项目已形成AI辅助的DRAM单粒子效应评估流程：限定自然语言请求映射到冻结
参数矩阵，真实Geant4电离沉积经等效源进入TCAD，再由匹配对照、参考支路、
数值检查和来源归档共同支持指标提取。
v1完成了小矩阵扫描及有限外推演示；v2与D1说明读出初态和参考选择会实质改变
结果解释；R50-main在独立物理修订后，完成参考、源隔离、代表点验证及
32个独立训练事件的64个配对响应，并冻结后续64条测试位预测。

四个训练岭模型在工况分组交叉验证中优于均值基准，说明基于沉积量、栅长和
方向的低复杂度代理值得继续验证。当前结论限于这一训练阶段和未校准二维模型；
域内留出、尺寸外推及工艺扩展须按原设计完成真实评分后再加入最终结论。
报告保留失败、未决与未完成状态，使后续更新可以增加证据，而不必改写历史。
课程最终交付还应补齐真实成员贡献、教师指导及作品演示记录。

\FloatBarrier
\clearpage
\bibliographystyle{unsrt}
\bibliography{references}
\clearpage
\appendix
\ctexset{section={name={附录,},number=\Alph{section}}}
\section{课程要求与现有证据对应}
\small
\begin{longtable}{p{0.21\textwidth}p{0.44\textwidth}p{0.27\textwidth}}
\toprule
课程要求&现有材料与报告位置&尚待补齐的证据\\
\midrule
\endfirsthead
\toprule
课程要求&现有材料与报告位置&尚待补齐的证据\\
\midrule
\endhead
理解DRAM、SEE与AI方法&第1至4章给出原理、简化假设、接口及回归方法&学生个人理解与答辩表现\\
2至3人合作&本报告提供可按模型、平台、数据分析划分的工作对象&真实成员、职责、完成日期与贡献\\
可行性、方案与进度&第1章、第4章和共享进度表；原设计与版本修订保留&教师对方案及变更的审核记录\\
软硬件及环境搭建&Windows与Linux VM的实际远程访问流程；Geant4、TCAD归档&实验室资源正式交接及使用说明\\
基础模型与参数来源&第2章参数表；自编几何和mixed-mode输入&真实DRAM模型或参考电学校准\\
自然语言参数遍历&v1请求到执行；R50的preview/audit/locate映射&学生现场操作与问题解释记录\\
Geant4到TCAD传递&式\eqref{cr:transfer}、源文件及积分验证&更细空间映射与材料栈适用性\\
多尺寸、粒子和方向评估&v1完整小矩阵；R50训练与冻结测试设计&R50其余测试及后续扩展最终验收\\
拟合、外推与验证&v1有限外推；R50四模型LOCO与测试前预测冻结&R50域内/域外完整真实评分\\
总结、演示与考核&本阶段报告及附录只读演示步骤&教师评审、答辩、成绩与验收事实\\
\bottomrule
\end{longtable}
\normalsize

\clearpage
\section{成员、指导与实际阶段进度记录}
原课程大纲列指导教师为李洋。本稿不填写未提供的学生姓名、学号、实际贡献、
审核日期或答辩结果；以下表格用于参与者依据真实工作补全。
AI辅助工具不计作项目小组成员。

\begin{table}[H]
\centering\small
\caption{真实课程信息登记表。空白表示尚未收到信息，不表示相应活动已经发生。}
\begin{tabularx}{\textwidth}{lXX}
\toprule
项目&真实记录&对应材料或日期\\
\midrule
成员1及学号&\rule{4cm}{0.4pt}&\rule{4cm}{0.4pt}\\
成员2及学号&\rule{4cm}{0.4pt}&\rule{4cm}{0.4pt}\\
成员3（如有）&\rule{4cm}{0.4pt}&\rule{4cm}{0.4pt}\\
各成员实际分工&\rule{4cm}{0.4pt}&\rule{4cm}{0.4pt}\\
指导教师&李洋（课程大纲所列）&实际指导记录待补\\
方案/变更审核&\rule{4cm}{0.4pt}&\rule{4cm}{0.4pt}\\
演示与答辩&\rule{4cm}{0.4pt}&\rule{4cm}{0.4pt}\\
\bottomrule
\end{tabularx}
\end{table}

\begin{table}[H]
\centering\small
\caption{实际阶段演进。日期由原始归档UTC时间换算为UTC$-08{:}00$；
当前完成数以正文共享快照表为准，run ID中的原UTC时间保留。}
\begin{tabularx}{\textwidth}{llX}
\toprule
本地日期&阶段&实际形成的主要成果\\
\midrule
2026-10-09&v1&16工况真实链路、有限预测外推、代表点组合细化及归档\\
2026-10-09&v2&受控初态、8独立事件与16配对响应、小样本读出判据\\
2026-10-09至10&D1&9项无粒子诊断、读出偏移和匹配参考实现证据\\
2026-10-10&100 $\Omega$/probe&保留预充失败；完成50 $\Omega$独立三项probe\\
2026-10-10至11&R50-main&参考、源隔离、数值检查、训练与测试前预测冻结\\
共享快照之后&待完成部分&按原冻结设计补齐测试评分、扩展验证和课程活动\\
\bottomrule
\end{tabularx}
\end{table}

\clearpage
\section{可重复演示步骤与材料入口}
演示从本报告证据包开始，不根据当前滚动运行状态临时换图或改数。
先展示模型参数表、版本和证据截止时间，再按下列顺序核对一个预登记事件。
所有只读命令均在仓库根目录运行；源码漂移或解析拒绝须保留原输出，不在现场
改条件凑成通过。

\begin{enumerate}
\setlength{\itemsep}{0.15em}
\item 打开本报告及共享\sourcepath{sources.json}，说明训练、预测与测试评分的区别，
并用正文进度表指出哪些部分仍未完成。
\item 运行R50自然语言映射预览和软件审计。预览展示64个独立事件及128个配对槽位；
审计验证限定模板和拒绝示例，不产生新deck或提交求解。
\item 定位预登记的$L=0.18\um$、$\alpha$、30°、首个事件的存储1响应，
从事件元数据核对seed、坐标、总电离沉积及三文件哈希。
\item 沿来源记录展示源项传递、匹配control与neutral，以及2.8 ns的$\Delta Q$
和4 ns的$M,\Delta M$。说明同一事件的0/1共享源，不增加独立事件数。
\item 展示训练LOCO表、模型及测试预测冻结记录。测试未完成时只展示已冻结预测的
身份和时间，不把预测图冒充真实测试点。
\item 展示历史D1参考比较和R50预充修订图，解释一个失败为什么必须保留，
以及改变物理版本为什么需要重新建立参考。
\end{enumerate}

\noindent 下列命令是可重复的只读演示入口，不表示学生现场演示已经发生：
{\small\begin{verbatim}
$py = 'local_runtime/vdmos-analysis-venv/Scripts/python.exe'
$demo = 'experiments/dram-see-v3/r50-main/demo.py'
& $py -B $demo preview --compact
& $py -B $demo audit
& $py -B $demo locate --event-label l018-alpha-a30-s202610311 --bit 1
& $py -B $demo preview --text "扫描几个尺寸"
\end{verbatim}}
最后一条预期拒绝，不能把非零退出码当成仿真失败。若需演示真实求解，
应另行确认唯一调度者、冻结输入和资源安排；本步骤不启动求解器。

\noindent 主要入口如下，均为仓库相对路径，具体内容摘要由共享来源清单绑定：
\begin{itemize}
\item 历史链路与预测：\sourcepath{experiments/dram-see-v1/delivery.md}。
\item 历史读出统计：\sourcepath{experiments/dram-see-v2/data/summary/delivery-index.json}。
\item D1诊断：\sourcepath{experiments/dram-readout-offset-d1/interpretation.md}。
\item R50冻结预测：\sourcepath{experiments/dram-see-v3/r50-main/data/predictions/freeze.json}。
\item 当前演示契约：\sourcepath{experiments/dram-see-v3/r50-main/demo-implementation.md}。
\item 本报告数据表、图件及来源：\sourcepath{reports/final/tables/}、
\sourcepath{reports/final/figures/}及\sourcepath{reports/final/sources.json}。
\end{itemize}

\end{document}
